\section{Zusammenfassung}
\label{sec:fazit}

In dieser Arbeit wurden zehn verschiedene Netzarchitekturen auf dem
PatternNet-Datensatz verglichen; die Experimente basierten dabei auf einer
verkleinerten, aber weiterhin balancierten Teilmenge.
Die Studie zeigt, dass architekturspezifische Entscheidungen (Tiefe,
Feature-Map-Anzahl, Normalisierung und Regularisierung) einen signifikanten
Einfluss auf Generalisierung und Stabilität des Trainings haben.
Die eingesetzte Datenaugmentation erhöht die Robustheit
gegenüber Lage- und Skalenvariationen und trägt zur Reduktion spezifischer
Verwechslungen bei. Die Variante \emph{BatchNorm\_Dense\_Pool\_Conv\_Dropout\_V2\_Net},
in der zusätzliche Layer und eine erhöhte Feature-Map-Anzahl verwendet wurden,
erforderte deutlich längere Trainingszeiten; die finale Bewertung ergab jedoch
keine relevante Verbesserung der Accuracy. Die Confusion-Matrizen zeigen, dass
die meisten Fehler auf visuell ähnliche Klassen zurückzuführen sind und nicht auf
ein systematisches globales Verwechslungsmuster.

Wesentliche Erkenntnisse:
\begin{itemize}
      \item Das beste Modell war \emph{BatchNorm\_Dense\_Pool\_Conv\_Net} (85.8\% Val, 84.2\% Test).
      \item Die komplexere Variante \emph{BatchNorm\_Dense\_Pool\_Conv\_Dropout\_V2\_Net} erforderte
            deutlich längere Trainingszeiten, brachte aber keine relevante Verbesserung der Accuracy.
      \item Die eingesetzte Datenaugmentation erhöht die Robustheit gegenüber Lage- und Skalenvariationen.
      \item Die Confusion-Matrizen zeigen, dass die verbleibenden Fehler überwiegend auf visuell
            ähnliche Klassen zurückzuführen sind.
\end{itemize}
